
This framework addresses a gap in ML research methodology: no formalized protocol exists for early viability assessment before resource commitment. Existing approaches test hypothesis variations, provide theoretical bounds, or optimize hyperparameters, but do not systematically gate non-viable hypotheses at micro-pilot scale.

\subsubsection{2.1 Ablation Studies}

Ablation studies are standard practice for testing hypothesis variations. A researcher might compare attention mechanisms with 1, 4, 8, or 16 heads to identify optimal configurations. These studies answer ''which variant performs best?'' not ''is this approach viable?'' The distinction is critical: ablations assume hypothesis viability and optimize within that space. Viability gates assess feasibility before committing to full implementation.

The ablation literature provides no formalized early-stop decision protocol. Computational infeasibility is typically discovered post-hoc after implementing all variants and measuring overhead across full-scale experiments. The h-e1 case (68.65\% overhead) exemplifies this pattern: the approach succeeded technically but failed on deployment constraints discoverable at micro-pilot scale.

\subsubsection{2.2 Computational Complexity Analysis}

Big-O notation characterizes asymptotic algorithmic scaling (O(n), O($n^{2})$, O(n log n)). While valuable for theoretical analysis, Big-O omits constant factors and implementation details that dominate practical performance. An O(n) algorithm with 68.65\% overhead constant may be less deployable than an O(n log n) algorithm with 5\% overhead.

Complexity analysis assumes idealized conditions: infinite memory, uniform data access, negligible cache effects. Real implementations encounter memory bottlenecks, I/O overhead, and hardware-specific factors that theoretical models do not capture. Empirical profiling addresses this gap by measuring actual overhead on target hardware. Current practice applies profiling post-implementation rather than at micro-pilot design stages.

\subsubsection{2.3 Early Stopping}

Early stopping literature primarily addresses training convergence: halting optimization when validation loss plateaus rather than continuing to a fixed epoch count. This optimizes hyperparameters (learning rate, batch size) to maximize accuracy while minimizing training cost.

The framework presented here applies early stopping to a different objective: viability assessment rather than accuracy optimization. The decision is binary (stop or continue based on threshold comparison) not continuous (find optimal hyperparameter value). The target is computational overhead rather than validation loss. While both paradigms share fail-fast principles, they address orthogonal problems: hyperparameter search assumes hypothesis viability and optimizes within constraints, while viability gates assess whether constraints can be met.

\subsubsection{2.4 Bayesian Optimization}

Bayesian optimization uses Gaussian processes to model expensive objective functions and sequentially select query points balancing exploration and exploitation. The framework has proven effective for hyperparameter tuning where each evaluation (training run) is costly.

The Bayesian update mechanism in the present framework shares sequential refinement principles but differs in application domain. Bayesian optimization maximizes an unknown function (find best hyperparameters). This framework predicts a known-in-principle quantity (full-scale overhead) from limited observations (micro-pilot measurements). Posterior refinement reduces prediction uncertainty rather than identifying optima. The gate structure (10 samples, 100 samples, full dataset) provides natural checkpoints for Bayesian updates.

\subsubsection{2.5 Positioning}

Pilot-Driven Viability Gates fills a methodological gap at the intersection of ablation studies, complexity analysis, and early stopping. The framework formalizes viability assessment preceding resource commitment, using empirical measurement to capture constant factors and hardware specifics missed by Big-O analysis, enabling incremental refinement through Bayesian updates, and providing binary stop/continue decisions based on threshold comparison with quantified confidence intervals.
